\documentclass{whiteboard}
\begin{document}
\begin{frame}[plain,t]
\bbcover{Programação Lógica}{Aritmética}{Prof. Edson Alves}{Campus UnB Gama: Faculdade de Ciências e Tecnologias em Engenharia}

\end{frame}
\begin{frame}[plain,t]
\begin{tikzpicture}
\node[draw,opacity=0] at (0, 0) {x};
\node[draw,opacity=0] at (14, 8) {x};

	\node[anchor=west] (title) at (0.0, 7.0) { \Large \bbbold{\textit{Backtracking}} };

\end{tikzpicture}
\end{frame}
\begin{frame}[plain,t]
\begin{tikzpicture}
\node[draw,opacity=0] at (0, 0) {x};
\node[draw,opacity=0] at (14, 8) {x};

	\node[anchor=west] (title) at (0.0, 7.0) { \Large \bbbold{\textit{Backtracking}} };


	\node[anchor=west] (a) at (1.0, 6.0) { $\star$ \bbtext{Quando Prolog tenta satisfazer um objetivo a respeito de um predicado, ele percorre} };

	\node[anchor=west] (a1) at (0.5, 5.5) { \bbtext{todas as cláusulas que definem o predicado} };

\end{tikzpicture}
\end{frame}
\begin{frame}[plain,t]
\begin{tikzpicture}
\node[draw,opacity=0] at (0, 0) {x};
\node[draw,opacity=0] at (14, 8) {x};

	\node[anchor=west] (title) at (0.0, 6.5) { \Large \bbbold{Referências} };

	\node[anchor=west] (a) at (1.0, 4.0) { $\star$ \bbbold{SWI-Prolog.} \bbenglish{https://www.swi-prolog.org/,} \bbtext{acesso em 10/03/2026.} };

	\node[anchor=west] (c) at (1.0, 3.0) { $\star$ \bbbold{Wikipédia.} \bbenglish{Prolog,} \bbtext{acesso em 10/11/2020.} };

	\node[anchor=west] (d) at (1.0, 5.0) { $\star$ \bbbold{MERRIT}\bbtext{, Dennis.} \bbenglish{Adventure in Prolog, Amzi!,} \bbtext{191 pgs, 2017.} };

\end{tikzpicture}
\end{frame}
\end{document}
